\chapter{Uvod}\label{cha:uvod}

Uvodno poglavje naj vsebuje vsebinski uvod, v katerem je pojasnjeno ozadje problema, in dve podpoglavji: \poglavjeref{sec:cilji_naloge}, \poglavjeref{sec:struktura_naloge}. Skupaj naj bo poglavje \poglavjeref{cha:uvod} dolgo vsaj 1,5 strani. Vsako novo poglavje naj se začne na lihi strani.

Uvod s predstavitvijo ozadja problema naj vsebuje vsaj en uvodni odstavek o obravnavani tematiki in predstavitvi izhodišča naloge.

\section{Cilji naloge}\label{sec:cilji_naloge}
Raziskovalne cilje predstavite v tem podpoglavju. Namesto ciljev so lahko zastavljene raziskovalne hipoteze (v tem primeru spremenite naslov poglavja na \ref{sec:cilji_naloge} \textbf{Raziskovalne hipoteze}).

V uvodu ne predstavljajte rezultatov in sklepov. V tem delu se osredotočite na to, kaj bo v delu obravnavano, kakšni so cilji/hipoteze. Pišite, kaj boste delali, kaj pričakujete od teoretičnih in kaj od praktičnih raziskav ter kakšna so tveganja in nevarnosti, da teh ciljev\\
ne boste dosegli.

\section{Struktura naloge}\label{sec:struktura_naloge}
V tem podpoglavju predstavite grobo strukturo dela, v kateri lahko na kratko z eno ali dvema povedma opišete vsebino posameznih poglavij. Spodaj je navedena struktura te predloge kot primer.

V tej predlogi je v nevsebinskem delu zaključne naloge (do strani xvi) potrebno izpolniti vse informacije, da bo naloga ustrezala podatkom kandidata oz. kandidatke ter podatkom o njegovi oz. njeni zaključni nalogi. Kazalo vsebine je po zaključku pisanja potrebno le osvežiti. Pri tem upoštevajte navodila glede uporabe ustreznih slogov, ki so navedena v \textbf{prilogi A}.

	V tej predlogi so \textbf{v vsebinskem delu} navodila in primeri za oblikovanje zaključne naloge, tehnične podrobnosti pa so dodatno opisane v \textbf{prilogi A}. Celotno besedilo te predloge mora študent oz.~študentka nadomestiti z besedilom, ki vsebinsko ustreza njegovi oz.~njeni zaključni nalogi. Naslove oštevilčenih poglavij na 1.~nivoju (\poglavjeref{cha:uvod}, \poglavjeref{cha:teorija} , \poglavjeref{cha:metodologija}, \poglavjeref{cha:rezultati}, \poglavjeref{cha:diskusija}, \poglavjeref{cha:zakljucki}, \textbf{\nameref{cha:literatura}}) je nujno potrebno pustiti v dokumentu. Poglavji \ref{cha:rezultati} in \ref{cha:diskusija} lahko združite v enotno poglavje \textbf{4 Rezultati in diskusija,} v kolikor je za izboljšanje jasnosti, razumevanja in berljivosti zaključnega dela to bolj primerno.